\documentclass[10pt,a4paper]{amsart}
\usepackage[margin=19mm]{geometry}
\usepackage{amsmath,amssymb,mathtools}
\usepackage[scr=rsfs,cal=euler]{mathalfa}
\usepackage{xfrac}
\usepackage[unicode,hidelinks,bookmarks=false]{hyperref}
\newcommand{\ie}{i.e.\ }
\newcommand{\cf}{cf.\ }
\newcommand{\resp}{respectively\ }
\newcommand{\Subsection}{\subsection}
\providecommand{\nobreakdash}{\nobreak\hbox{-}}
\setlength{\emergencystretch}{2em}
\multlinegap0pt
\usepackage{amssymb}
\usepackage{bm}
\newtheorem{lem}{Lemma}
\newcommand{\ev}{\operatorname{ev}}
\title{Rational higher order Whitehead products of mapping spaces and defining systems}
\author{Takahito Naito}
\date{Source: arXiv:2607.27663v1 (CC BY 4.0)}
\begin{document}
\maketitle

\noindent\emph{Source and licence.} Rational higher order Whitehead products of mapping spaces and defining systems. Version arXiv:2607.27663v1 (CC BY 4.0), distributed by arXiv under the Creative Commons Attribution 4.0 International licence, http://creativecommons.org/licenses/by/4.0/, as recorded in the arXiv licence metadata for this exact version and shown on the abstract page https://arxiv.org/abs/2607.27663v1. Full licence notice: \texttt{licenses/arxiv-2607.27663-CC-BY-4.0.txt}.\par\smallskip

\noindent\textit{Editorial scope.} Counted unit: the article's lem lem:strict-cohomology-model (source0.tex lines 1190--1204). The article's complete original proof of it is delivered with it (source0.tex lines 1206--1375, one \texttt{proof} environment of 170 lines). No mathematics is written, repaired or re-derived: the statement and the proof are the article's own lines, in the article's own notation, and no retained line is substituted or re-flowed.  No further source-local context is needed: every cross-reference of every retained line resolves inside the two delivered spans.  Nothing was omitted from any retained span, so the package carries no omission record.  No retained line contains a citation, so the wrapper carries no bibliography and no bibliography entry.  No retained line contains a figure environment, an image-inclusion command or drawing code, so no image is delivered and the package declares no asset dependency.  Editorial formatting only, in addition: an AMS \texttt{amsart} preamble that restates the article's own declarations and macros used by the retained lines, namely the environments \texttt{lem} on separate counters, together with the standard packages those lines need.  The retained environments print in this document as Lemma 1 in order of appearance, whereas the article prints the same items with the numbering of its own full document.  The complete native source of the article as distributed in the arXiv e-print for this exact version is bundled unchanged as \texttt{sources/206362/source0.tex}.
\begin{lem}\label{lem:strict-cohomology-model}
Let $\varphi':\wedge V\to \mathcal H_T\otimes A$ be a CDGA morphism.
Suppose that $(\rho\otimes 1)\varphi'$ is homotopic to a CDGA morphism
$
\psi:\wedge V\to \mathcal H_W\otimes A.
$
Then there exists a CDGA morphism
$
\varphi:\wedge V\to \mathcal H_T\otimes A
$
which is homotopic to $\varphi'$ and satisfies
$
(\rho\otimes 1)\varphi=\psi.
$
\end{lem}

\begin{proof}
Choose a CDGA homotopy
\[
H':
\wedge V
\longrightarrow
\mathcal H_W\otimes A\otimes \Lambda(t,dt)
\]
from $(\rho\otimes 1)\varphi'$ to $\psi$. Thus
$
\operatorname{ev}_0H'=(\rho\otimes 1)\varphi'
$
and
$
\ev_1H'=\psi.
$
Here $\ev_i:\Lambda(t,dt)\to \mathbb Q$ denotes the CDGA map given by $\operatorname{ev}_i (t) = i$.

It is enough to construct a CDGA homotopy
\[
H:
\wedge V
\longrightarrow
\mathcal H_T\otimes A\otimes \Lambda(t,dt)
\]
such that
$
\operatorname{ev}_0H=\varphi'
$
and
$
(\rho\otimes 1\otimes 1)H=H'.
$
The homotopy $H$ will be defined on the generators as follows.
For $v\in V$, write
\[
H'(v)
=
1\otimes H'_\emptyset(v)
+
\sum_{i=1}^p e_i\otimes H'_i(v).
\]
Define the components of $H(v)$ by
\[
H(v)
=
1\otimes H_\emptyset(v)
+
\sum_{i=1}^p e_i\otimes H_i(v)
+
\sum_{\substack{I\subsetneq[p]\\ |I|\geq 2}}
e_I\otimes H_I(v),
\]
where
$
H_\emptyset(v)=H'_\emptyset(v)
$
and
$
H_i(v)=H'_i(v)
$
for $i=1,\ldots,p$. It remains to define $H_I(v)$ for $|I|\geq 2$.
The construction proceeds by induction over an ordered basis of $V$ compatible
with the Sullivan filtration. Suppose first that $dv=0$.
Write
\[
        \varphi'(v)=\sum_{I\subsetneq[p]} e_I\otimes \varphi'_I(v).
\]
Then each $\varphi'_I(v)$ is a cocycle. In this case, for $|I|\geq 2$, define
\[
        H_I(v)=\varphi'_I(v)\otimes 1.
\]
Now suppose that $dv\neq 0$. Let $V_{<v}$ be the subspace spanned by the
generators preceding $v$.
By induction, assume that $H$ has been constructed
as a CDGA morphism on $\wedge V_{<v}$ and satisfies
$
\ev_0H=\varphi'
$
and
$
(\rho\otimes 1\otimes 1)H=H'
$
on $\wedge V_{<v}$. Since $dv\in\wedge V_{<v}$, the element $H(dv)$ is
already defined.
For $I\subsetneq[p]$ with $|I|\geq 2$, let $H_I(dv)$ be the
coefficient of $e_I$ in $H(dv)$, and put
\[
        \alpha_I=(-1)^{|e_I|}H_I(dv).
\]
Since $H$ is a CDGA morphism on $\wedge V_{<v}$,
$
        DH(dv)=H(d^2v)=0,
$
where $D$ denotes the differential of $A\otimes\Lambda(t,dt)$. 
Thus $\alpha_I$ is a cocycle in $A\otimes\Lambda(t,dt)$. 
Moreover, the induction hypothesis gives
\[
\operatorname{ev}_0(\alpha_I)
=
(-1)^{|e_I|}\varphi'_I(dv)
=
d_A\varphi'_I(v),
\]
where the last equality follows because $\varphi'$ is a CDGA morphism.

We use the standard
homotopy operator
$
K:A\otimes\Lambda(t,dt)\to A\otimes\Lambda(t,dt)
$
of degree $-1$ associated with evaluation at $t=0$. Thus
\[
        DK+KD=\operatorname{id}-\iota\operatorname{ev}_0
        \quad
        \text{and}
        \quad
        \operatorname{ev}_0K=0,
\]
where $\iota:A\to A\otimes\Lambda(t,dt)$ is given by $a\mapsto a\otimes 1$.
Here $K$ is used only as a linear map on the underlying cochain complex.
Put
\[
        \beta_I=\alpha_I-D(\varphi'_I(v)\otimes 1).
\]
Then $\beta_I$ is a cocycle and $\ev_0(\beta_I)=0$. Hence
$DK(\beta_I)=\beta_I$. Define
\[
H_I(v)
=
\varphi'_I(v)\otimes 1
+
K(\beta_I).
\]
Since $\ev_0K=0$, we have
$
\operatorname{ev}_0H_I(v)=\varphi'_I(v)
$
and
\[
        DH_I(v)
        =
        D(\varphi'_I(v)\otimes 1)+\beta_I
        =
        \alpha_I
        =
        (-1)^{|e_I|}H_I(dv).
\]
Since $\wedge V$ is free, this choice of $H(v)$ extends the partially defined
map uniquely to a graded algebra map.
The identity
$DH(v)=H(dv)$ makes this extension a CDGA morphism.
Continuing the
induction gives the required CDGA homotopy $H$.
Finally, set
\[
        \varphi=\ev_1 H.
\]
Then $\varphi$ is homotopic to $\varphi'$, and
\[
(\rho\otimes 1)\varphi
=
\ev_1 \bigl((\rho\otimes 1\otimes 1)H\bigr)
=
\ev_1 H'
=
\psi.
\]
This completes the proof.
\end{proof}

\end{document}
